\documentclass[ignorenonframetext,]{beamer}

%Set up notes
\setbeamertemplate{note page}[plain]
\setbeameroption{hide notes}

\usepackage{amssymb,amsmath}
\usepackage{ifxetex,ifluatex}
\usepackage{fixltx2e} % provides \textsubscript
\usepackage{tikz}
\usepackage{tikz-qtree}
\usepackage{pdfpages} %include beamer pdf slides from other presentations
\usepackage{color}
\ifxetex \usepackage{fontspec,xltxtra,xunicode}
  \defaultfontfeatures{Mapping=tex-text,Scale=MatchLowercase}
\else
  \ifluatex
    \usepackage{fontspec}
    \defaultfontfeatures{Mapping=tex-text,Scale=MatchLowercase}
  \else
    \usepackage[utf8]{inputenc}
  \fi
\fi

% Comment these out if you don't want a slide with just the
% part/section/subsection/subsubsection title:
\AtBeginPart{
  \let\insertpartnumber\relax
  \let\partname\relax
  \frame{\partpage}
}
\AtBeginSection{
  \let\insertsectionnumber\relax
  \let\sectionname\relax
  \frame{\sectionpage}
}
\AtBeginSubsection{
  \let\insertsubsectionnumber\relax
  \let\subsectionname\relax
  \frame{\subsectionpage}
}

\setlength{\parindent}{0pt}
\setlength{\parskip}{6pt plus 2pt minus 1pt}
\setlength{\emergencystretch}{3em}  % prevent overfull lines
\setcounter{secnumdepth}{0}

\title{Putting theory to work: The standard model}
\author{Instructor: David Jinkins}
\date{Date: Sept. 16, 2014}

\begin{document}

\frame{\titlepage}

\begin{frame}
\begin{itemize}
\itemsep1pt\parskip0pt\parsep0pt
\item
  Last time: The Heckscher-Ohlin Model
  \begin{itemize}
        \item Introduction
        \item The point
        \item The model
        \begin{itemize}
            \item Setup and assumptions
            \item Production possibilities
            \item Production and prices
            \item Autarchy equilibrium
            \item Trade equilibrium
        \end{itemize}
        \item The big four theorems
        \item The problems 
        \begin{itemize}
            \item Assumptions
            \item Evidence
        \end{itemize}
        \item Moving forward
    \end{itemize}
\end{itemize}

\end{frame}

\begin{frame}
\begin{itemize}
\itemsep1pt\parskip0pt\parsep0pt
\item
  Today: The Standard Model 
  \begin{itemize}
        \item Introduction
        \item Model
        \begin{itemize}
            \item Setup and assumptions
            \item Production possiblities 
            \item Supply and Prices
            \item Demand and Prices
            \item Trade and Prices
        \end{itemize}
        \item Growth
        \begin{itemize}
            \item Production
            \item Terms of Trade
            \item International effects
        \end{itemize}
        \item Tariffs and subsidies
        \item International lending
    \end{itemize}
    \item The location of production
\end{itemize}

\end{frame}

\begin{frame}
    \begin{itemize}
        \item But first a review!
    \end{itemize}
\end{frame}


\frame[plain]{\includegraphics[page=28,width=\textwidth]{heckscher_ohlin.pdf}}
\frame[plain]{\includegraphics[page=29,width=\textwidth]{heckscher_ohlin.pdf}}
\frame[plain]{\includegraphics[page=30,width=\textwidth]{heckscher_ohlin.pdf}}
\frame[plain]{\includegraphics[page=31,width=\textwidth]{heckscher_ohlin.pdf}}
\frame[plain]{\includegraphics[page=32,width=\textwidth]{heckscher_ohlin.pdf}}
\frame[plain]{\includegraphics[page=34,width=\textwidth]{heckscher_ohlin.pdf}}
\frame[plain]{\includegraphics[page=35,width=\textwidth]{heckscher_ohlin.pdf}}
\frame[plain]{\includegraphics[page=38,width=\textwidth]{heckscher_ohlin.pdf}}
\frame[plain]{\includegraphics[page=39,width=\textwidth]{heckscher_ohlin.pdf}}
\frame[plain]{\includegraphics[page=44,width=\textwidth]{heckscher_ohlin.pdf}}
\frame[plain]{\includegraphics[page=45,width=\textwidth]{heckscher_ohlin.pdf}}
\frame[plain]{\includegraphics[page=46,width=\textwidth]{heckscher_ohlin.pdf}}
\frame[plain]{\includegraphics[page=48,width=\textwidth]{heckscher_ohlin.pdf}}
\frame[plain]{\includegraphics[page=49,width=\textwidth]{heckscher_ohlin.pdf}}
\frame[plain]{\includegraphics[page=53,width=\textwidth]{heckscher_ohlin.pdf}}
\frame[plain]{\includegraphics[page=54,width=\textwidth]{heckscher_ohlin.pdf}}
\frame[plain]{\includegraphics[page=55,width=\textwidth]{heckscher_ohlin.pdf}}
\frame[plain]{\includegraphics[page=56,width=\textwidth]{heckscher_ohlin.pdf}}
\frame[plain]{\includegraphics[page=67,width=\textwidth]{heckscher_ohlin.pdf}}

\begin{frame}
    \begin{itemize}
        \item End review!
    \end{itemize}
\end{frame}

\begin{frame}{Chapter 6: The standard model}

    \begin{itemize}
        \item Best chapter yet 
        \begin{itemize}
            \item Probably a Krugman chapter
            \item Lots of intuition, simple model, important questions
        \end{itemize}
        \item Why do countries trade?
        \begin{itemize}
            \item Differences in technology?
            \item Differences in factor endowments?
        \end{itemize}
        \item We just need different production possibilities
        \item For much analysis \emph{it doesn't matter!}
    \end{itemize}
            
\end{frame}

\begin{frame}{Laying out the model}

    \begin{itemize}
        \item Environment 
        \item Supply and Prices
        \item Demand, Welfare, and Prices
        \item Equilibrium
    \end{itemize}
            
\end{frame}

\begin{frame}{Environment}

    \begin{itemize}
        \item Two countries: Home (H), Foreign (G)
        \item Two goods: Food (F), Clothes (C)
        \item We aren't going to specify the factors
        \begin{itemize}
            \item Ricardo: Labor
            \item Specific Factors: Land, Capital, Labor
            \item Heckscher Ohlin: Capital, Labor
        \end{itemize}
        \item We aren't going to specify exact technology for converting factors into goods
        \begin{itemize}
            \item Ricardo: Unit Labor requirement
            \item Specific Factors: $Q_C(K,L), Q_F(T,L)$
            \item Heckscher Ohlin: $Q_C(K,L), Q_F(K,L)$
        \end{itemize}

    \end{itemize}
            
\end{frame}

\begin{frame}{Supply}

    \begin{itemize}
        \item Each country is going to be given a production possibilities frontier
        \item We don't ask where it came from
        \item Book: Economy will produce at point on PPF which maximizes value
        \item This result makes sense for a dictator (or social planner)
        \item Will the market deliver that result?
    \end{itemize}

\end{frame}

\begin{frame}{Supply}

    \begin{itemize}
        \item We know: 
        \begin{enumerate}
            \item Full employment of factors (and positive marginal product) means production on the PPF 
            \item We have argued that slope of PPF is the negative ratio of marginal products
            \item From wages/rental, in equlibrium: $\frac{P_F}{P_C} = \frac{MPL_C}{MPL_F}$
        \end{enumerate}
        \item Therefore: a line with negative slope of prices will be tangent to the PPF at the point of production
        \item But what about the value of production?
    \end{itemize}

\end{frame}

\begin{frame}{The isovalue lines}

    \begin{itemize}
        \item What is the equation for the price line tangent to the PPF?
    \end{itemize}

\end{frame}

\begin{frame}{Supply}

    \begin{itemize}
        \item We have shown that equilibrium production is at a point where production value is maximized
        \item Follows from the equilibrium wage and rental equations, and therefore the no-profit condition
    \end{itemize}

\end{frame}

\begin{frame}{Supply}

    \includegraphics[scale=0.25]{isovalue.png}

\end{frame}

\begin{frame}{Supply}

    \includegraphics[scale=0.25]{price_supply.png}

\end{frame}

\begin{frame}{Supply}

    \includegraphics[scale=0.25]{rprq.png}

\end{frame}

\begin{frame}{Demand}

    \begin{itemize}
        \item Until this class, demand has just been a curve
        \item (ok, in a homework I made you plot a couple of them)
        \item Today we want to analyze welfare 
        \item We need to go a bit farther
        \item In short -- today we skimp on the supply, but say more about demand
    \end{itemize}

\end{frame}

\begin{frame}{Demand}

    \begin{itemize}
        \item We have just seen that the value produced is $P_C Q_C + P_F Q_F$.
        \item Trade balance says that the value of consumption must equal value of production:
        \begin{equation*}
            P_C Q_C + P_F Q_F = P_C D_C + P_F D_F
        \end{equation*}
        \item What does this imply about the location of consumption on our PPF when trading?
        \item What additional condition do we have in autarchy?
    \end{itemize}

\end{frame}

\begin{frame}{The indifference curve}

    \begin{itemize}
        \item The set of bundles of goods over which the consumer is indifferent 
        \begin{itemize}
        \item More of either good is better
        \item Slope at any point in $(-\infty,0]$, that is, must compensate
        \item If I have little of one good, you have to give me a lot of the other to compensate for further reduction. (decreasing marginal utility)
        \end{itemize}
        \item Results in convex curves emanating from the origin
    \end{itemize}

\end{frame}

\begin{frame}{Demand}
    \begin{itemize}
        \item Where would autarchy consumption be?
    \end{itemize}
    \includegraphics[scale=0.25]{ind_curves.png}

\end{frame}

\begin{frame}{Price changes and welfare}

    \begin{itemize}
        \item Suppose my country sells clothes to buy food
        \item If the relative price of clothes goes up:
        \begin{itemize}
            \item I become richer -- Foreign has to give me more food for my clothes (\emph{income effect})
            \item I change my consumption bundle to relatively more food (\emph{income effect})
            \item What happens to production at Home?
            \item What happens to production at Foreign?
        \end{itemize}
    \end{itemize}

\end{frame}

\begin{frame}{Price changes and welfare}

    \begin{itemize}
        \item See income and substitution effects in chart
    \end{itemize}
    \includegraphics[scale=0.25]{price_change_demand.png}

\end{frame}

\begin{frame}{Price changes and welfare}

    \includegraphics[scale=0.25]{rs_rd.png}
    \begin{itemize}
        \item Which good does Home export?
    \end{itemize}

\end{frame}

\begin{frame}{Price changes and welfare}

    \begin{itemize}
        \item How would Home consumption change if the relative price of Clothes fell?
    \end{itemize}

\end{frame}

\begin{frame}{Price changes and welfare}

    \begin{itemize}
        \item \emph{Terms of trade}
        \item Defined as the price of the exported good divided by the price of the imported good
        \item Welfare increases if terms of trade rise (or improve)
        \item Welfare decreases if terms of trade fall (or worsen)
        \item A worsening of the terms of trade can never make welfare fall below autarchy level (why?)
    \end{itemize}

\end{frame}

\begin{frame}{Equilibrium Trade}

    \includegraphics[scale=0.25]{rs_world.png}
    \begin{itemize}
        \item Are $RS$ and $RS^*$ the autarchy supplies?
    \end{itemize}

\end{frame}

\begin{frame}{Equilibrium Trade}

    \includegraphics[scale=0.25]{trade_flows.png}
    \begin{itemize}
        \item What is wrong with these pictures (under the assumptions we have typically been making?)
    \end{itemize}

\end{frame}

\begin{frame}{End model development}

    \begin{itemize}
        \item We have now presented the model
        \item Simpler than the others we developed 
        \item Countries have PPFs
        \item Countries have demand
        \item Prices affect both production and consumption 
        \item Improvement in terms of trade helps
        \item Worsening of terms of trade hurts
    \end{itemize}

\end{frame}

\begin{frame}{Some analysis}

    \begin{itemize}
        \item Is growth in other countries good or bad for us?
        \item If other countries will not respond, should we impose tariffs?
        \item How will Danish current account surplus affect future Danish consumption? 
    \end{itemize}

\end{frame}

\begin{frame}{Growth}

    \begin{itemize}
        \item Question 1: Is China's growth good for Denmark?
        \item Question 2: Is Danish growth better or worse thanks to international trade?
    \end{itemize}

\end{frame}

\begin{frame}{Growth} 

    \begin{itemize}
        \item Claim in book (pg. 158):

            \emph{The international trade effects of growth result from the fact that such growth typically has a bias.}

        \item I suspect that unbiased growth still affects trade
        \item come back to this point
    \end{itemize}

\end{frame}

\begin{frame}{Biased growth}

    \begin{itemize}
        \item Growth that expands the production possibilities more towards one good
        \item Growth that, for any fixed price, always causes relatively more of one good to be produced
    \end{itemize}

\end{frame}

\begin{frame}{Biased growth}

    \includegraphics[scale=0.25]{biased_growth.png}

\end{frame}

\begin{frame}{Biased growth}

    \begin{itemize}
        \item Suppose equilibrium trade prices didn't change
        \item Foreign supplies same good mix
        \item Home supplies relatively more of one good 
        \item This will affect the world relative supply 
    \end{itemize}

\end{frame}

\begin{frame}{Biased growth}

    \includegraphics[scale=0.25]{biased_growth_rs.png}

\end{frame}

\begin{frame}{Biased growth}

    \begin{itemize}
        \item Of course, going to affect equilibrium price
    \end{itemize}
    \includegraphics[scale=0.25]{biased_growth_rs.png}

\end{frame}

\begin{frame}{Biased growth}

    \includegraphics[scale=0.25]{cloth_growth.png}

\end{frame}

\begin{frame}{Biased growth}

    \includegraphics[scale=0.25]{food_growth.png}

\end{frame}

\begin{frame}{Biased growth}

    \begin{itemize}
        \item Effect on price depends upon which good gets the biased growth
        \item Effect on price does not depend upon which country grows
        \item If Home exports cloth, does growth in Foreign in Food help or hurt Home? Why?
    \end{itemize}

\end{frame}

\begin{frame}{Biased growth}

    \begin{itemize}
        \item Export-biased growth is bad for terms of trade
        \item Import-biased growth is good for terms of trade
    \end{itemize}

\end{frame}

\begin{frame}{Biased growth}

    \begin{itemize}
        \item Columbia Economist Jagdish Baghwati showed as a graduate student that growth at home can actually reduce welfare
        \item But productivity growth usually increases income, even if it worsens terms of trade
        \item Growth of other countries in Home's export industry does not increase income 
    \end{itemize}

\end{frame}

\begin{frame}{Unbiased Growth}

    \begin{itemize}
        \item I think it also affects trade and production
        \item Effect similar to export-orientated biased growth (bad)
    \end{itemize}

\end{frame}

\begin{frame}{Case study from book}

    \begin{itemize}
        \item Should we expect growth in developing countries to hurt developed world consumers?
        \item Depends in which industry their growth happens
        \item Case: Chinese productivity growth in manufacturing -- China's export industry
    \end{itemize}

\end{frame}

\begin{frame}{Case study from book}

    \begin{itemize}
        \item China-US terms of trade
    \end{itemize}
    \includegraphics[scale=0.20]{china_us_tot.png}
    \begin{itemize}
        \item What the book isn't saying\dots
    \end{itemize}

\end{frame}

\begin{frame}{Tariffs and Export Subsidies}

    \begin{itemize}
        \item What are they?
        \begin{itemize}
            \item \emph{Tariffs} are taxes levied on imports
            \item \emph{Export subsidies} are payments given to domestic producers who sell abroad
        \end{itemize}
        \item Scope - effect on terms of trade
        \begin{itemize}
            \item Tariffs are not usually motivated by terms of trade
            \item They do have terms of trade effects
        \end{itemize}
        \item Questions
        \begin{itemize}
            \item Do tariffs help or hurt Home consumers?
            \item Do export subsidies help or hurt Home consumers?
        \end{itemize}
    \end{itemize}

\end{frame}

\begin{frame}{Tariffs}
    
    \begin{itemize}
        \item Tariffs are a tax applied in some places, not others
        \item Different relative prices faced by people in different countries
        \item We defined terms of trade as export good pr. / import good pr.
        \item But now we have two sets of rel. prices -- \emph{internal} and \emph{external}
    \end{itemize}

\end{frame}

\begin{frame}{Tariffs}
    
    \begin{itemize}
        \item Tariffs are a tax applied in some places, not others
        \item Different relative prices faced by people in different countries
        \item We defined terms of trade as export good pr. / import good pr.
        \item But now we have two sets of rel. prices -- \emph{internal} and \emph{external}
    \end{itemize}

\end{frame}

\begin{frame}{Tariffs}
    
    \begin{itemize}
        \item We are going to call \emph{external} prices the terms of trade
        \begin{itemize}
            \item External prices govern the value of Home's exports in terms of Foreign's exports
            \item Internal prices govern the choices of production and consumption at Home
        \end{itemize}
        \item Suppose Home puts a tariff on Food
        \item Suppose external prices do not change
        \begin{itemize}
            \item What happends to production of Food at home?
            \item What happends to consumption of Food at home?
            \item If Food were Labor-intensive, which factor benefits?
        \end{itemize}
    \end{itemize}

\end{frame}

\begin{frame}{Tariffs}
    
    \begin{itemize}
        \item We are going to call \emph{external} prices the terms of trade
        \begin{itemize}
            \item External prices govern the value of Home's exports in terms of Foreign's exports
            \item Internal prices govern the choices of production and consumption at Home
        \end{itemize}
        \item Suppose Home, Clothes exporter, puts a tariff on Food
        \item Suppose external prices do not change
        \begin{itemize}
            \item What happends to production of Food at home?
            \item What happends to consumption of Food at home?
            \item If Food were Labor-intensive, which factor benefits?
        \end{itemize}
    \end{itemize}

\end{frame}

\begin{frame}{External Price (TOT) Effect of Home Tariff on Food}

    \includegraphics[scale=0.25]{tariff_tot.png}

\end{frame}

\begin{frame}{External Price (TOT) Effect of Home Tariff on Food}

    \begin{itemize}
        \item Terms of Trade for Home improved
        \item Now Home can afford more Food for the same production of Clothes
        \item Great, right! But wait, let's think about this again
        \item Domestically, a Tariff is initially like a worsening of terms of trade!
        \item We need the External Price TOT Effect to dominate 
        \item There is an optimal tariff
        \begin{itemize}
            \item Typically small for small countries, large for large countries.  Why?
            \item Talk more about this in a later chapter
        \end{itemize}
    \end{itemize}

\end{frame}

\begin{frame}{Export subsidy}

    \begin{itemize}
        \item In a similar fashion, will reduce external TOT
        \item Domestically, however, will increase the internal price of the export good
    \end{itemize}

    \includegraphics[scale=0.20]{export_subsidy_tot.png}

\end{frame}

\begin{frame}{Trade Policy Discussion}

    \begin{itemize}
        \item Export subsidy and tariff seem very similar
        \item That is because we are missing something really important \end{itemize}

\end{frame}

\begin{frame}{Trade Policy Discussion}

    \begin{itemize}
        \item Financing! 
        \begin{itemize}
            \item Tariffs are financed by Foreign
            \item Export Subsidies are financed by Home
        \end{itemize}
    \end{itemize}

\end{frame}

\begin{frame}{Trade Policy Discussion}

    \begin{itemize}
        \item Is it good for Denmark if China subsidizes its exports?
        \item The (Stiglitz, I think) new car example
        \item Good for who again?
    \end{itemize}

\end{frame}

\begin{frame}{Intertemporal Trade}

    \begin{itemize}
        \item Dipping into 2nd half of the course -- International Finance
        \item Why? Because analysis is similar to current model
        \item Setup:
        \begin{itemize}
            \item Two countries: Home (H), Foreign (G)
            \item Two goods: Cheese today, and Cheese tomorrow
            \item Countries have different technology for converting today Cheese into tomorrow Cheese.
            \item Maybe only Home has refrigerators
        \end{itemize}
    \end{itemize}
            
\end{frame}

\begin{frame}{Intertemporal Trade}

    \begin{itemize}
        \item More serious motivation:
        \begin{itemize}
            \item One country might have better investment opportunities
            \item One country may have found oil, but it will run out tomorrow
        \end{itemize}
    \end{itemize}
            
\end{frame}

\begin{frame}{Intertemporal PPF}

    \includegraphics[scale=0.25]{intertemporal_ppf.png}

\end{frame}

\begin{frame}{Price of Future Consumption}

    \begin{itemize}
        \item Just like bank loan, you have to pay for current consumption 
        \item Suppose that you have to pay $1+r$ units of future cheese for one unit of current cheese
        \item How much current cheese to I have to give up for one unit of future cheese?
        \item That is the relative price of future cheese in terms of current cheese
    \end{itemize}

\end{frame}

\begin{frame}{Intertemporal Cheese equilibrium}

    \includegraphics[scale=0.25]{intertemporal_cheese.png}

\end{frame}

\begin{frame}
\begin{itemize}
\itemsep1pt\parskip0pt\parsep0pt
\item
  Summary 
  \begin{itemize}
        \item We set up the model
        \begin{itemize}
            \item More on demand
            \item Less on supply
        \end{itemize}
        \item Growth
        \item Tariffs and export subsidies
        \item International cheese lending
    \end{itemize}
\end{itemize}

\end{frame}

\begin{frame}
\begin{itemize}
\itemsep1pt\parskip0pt\parsep0pt
\item
  Next time: 
  \begin{itemize}
        \item The location of production
        \begin{itemize}
            \item One of the things Krugman won Nobel for 
            \item Getting behind on the syllabus (I'll update it)
        \end{itemize}
        \item Monopolistic Competition
        \begin{itemize}
            \item So far firms have not made profits
            \item Now they will, but still not monopolies
            \item Sometimes called imperfect competition
            \item Will 'explain' two-way trade puzzle
        \end{itemize}
    \end{itemize}
\end{itemize}

\end{frame}

\end{document}






